\section{Introduction}
\label{sec:intro}

The Internet of Things now connects tens of billions of constrained
devices, and every gateway that bridges those devices to the wider
network is a high-value attack surface. Flow-level intrusion detection
at the gateway is the standard line of defence, and recent benchmarks
report near-perfect closed-set accuracy for gradient-trained
classifiers on IoT traffic. The deployment story is less settled.
A modern IoT gateway runs on a few hundred megabytes of memory and a
single-digit-watt power envelope, and the same gradient-trained
classifier that scores well in the benchmark must coexist with the
gateway's primary routing workload at line rate. Large language models
have transformed reasoning over heterogeneous tabular inputs and would
in principle generalise across attack families without per-dataset
retraining, yet a foundation-model inference per packet is
incompatible with any realistic IoT edge budget.

The Integrated Sensing, Memory, Communication, and Computation (SMCC)
paradigm provides the architectural lens that resolves this tension.
SMCC separates the offline ingestion and compression of training
experience, the \emph{memory} stage, from the online line-rate
response, the \emph{computation} stage. A detector designed under
SMCC need not run the same machinery at training time and at
inference time. Expensive reasoning happens once, offline, against a
compact compressed view of the training distribution. The deployed
component carries no model weights and issues no foundation-model
call per flow. This split is what makes large-AI deployment on
resource-constrained IoT gateways tractable rather than aspirational.

We instantiate this split for IoT intrusion detection.
A retrieval-augmented language model acts as the memory stage.
\texttt{claude-haiku-4-5-20251001} (Claude Haiku 4.5) reads
class-representative neighbours retrieved from a Chroma vector store
of training flows and synthesises interpretable threshold rules once,
offline, as a JSON policy. A stateless rule evaluator acts as the
computation stage. It applies the policy to live flows as vectorised
threshold comparisons that map directly onto access-control list
firewall entries. No language-model call is required at inference
time, and the runtime dependency is a few dozen comparisons per flow.
The full offline policy-generation cost is approximately \$0.05 per
dataset at current Claude Haiku 4.5 API rates, which is the price of
a single training run rather than a recurring inference budget.

The architecture is viable because retrieval-augmented prompts can
recover the discriminative threshold structure of a benchmark IoT
dataset from a small representative sample without gradient training.
On datasets with well-separated traffic signatures the resulting
policy matches Random Forest within $0.06$ macro-F1 while running
4.5$\times$ to 80$\times$ faster per flow. On datasets with high
intra-class overlap the policy degrades gracefully, and the same
retrieval signal supplies a confidence score that flags
out-of-distribution traffic. We name this confidence score
Suspicious Multi-Class Confidence (SMCC, by construction aligned with
the special-issue acronym), and use it to extend the closed-set
detector to zero-day flows without retraining.

The paper makes four contributions.

\begin{itemize}
  \item An SMCC-aligned architecture for IoT intrusion detection that
    confines all foundation-model computation to an offline memory
    stage and deploys the result as a stateless threshold-rule
    evaluator. The runtime evaluator runs at line rate on commodity
    gateway hardware and requires no model weights.

  \item A multi-class evaluation across five IoT benchmarks
    (CIC-IoT2023, WUSTL-IIoT, TON\_IoT, Bot-IoT, UNSW-NB15) that
    characterises the operating envelope of LLM-generated threshold
    rules. The LLM matches Random Forest within $0.06$ macro-F1 on
    four of five datasets and exposes a structural expressiveness
    ceiling on the fifth.

  \item A retrieval-confidence protocol that extends the closed-set
    detector to zero-day attack families. The Suspicious Multi-Class
    Confidence score isolates the withheld attack class with
    $\text{ZDR} \geq 0.999$ on two of five withholding scenarios
    without any retraining of the policy.

  \item A feature-name anonymization ablation and a Kerckhoffs-style
    adversarial robustness study that together address the
    trustworthiness axis of SMCC. Anonymizing feature names does not
    degrade and in three of five cases improves detection F1, and
    LLM-generated rules distribute evasion cost across multiple
    features rather than concentrating it in a single bottleneck.
\end{itemize}

The remainder of the paper is organised as follows.
Section~\ref{sec:related} surveys the closest prior work on
LLM-assisted intrusion detection, retrieval-augmented generation, and
edge deployment of foundation models.
Section~\ref{sec:prelim} fixes notation and frames the SMCC
memory-computation split that the rest of the paper exploits.
Section~\ref{sec:architecture} describes the offline policy-synthesis
pipeline and the online rule evaluator that together realise that
split. Section~\ref{sec:framework} formalises the four evaluation
protocols. Section~\ref{sec:resultsAndDisc} reports results across
five IoT benchmarks and analyses the conditions under which the
approach is deployable. The appendix contains per-class breakdowns
and additional ablation tables.
